% GENERATED, edit the French unit or the translation table
\mxnewpage\section*{EV2-A -- Mock exam 1}\label{EV2-A}
\addcontentsline{toc}{section}{EV2-A -- Mock exam 1}
\textbf{Paper reserved for the teacher}. Chapters 1 to 7.\par Duration: 2 h. No calculator. Marked out of 20. 14--16 January 2027.
\subsection*{Part 1 -- Fluency (6 points)}
Twelve questions, 0.5 point each. Only one correct answer; no justification required. An incorrect or missing answer scores 0, with no loss of points.
\begin{enumerate}
\item Calculate $4+2\times3$.\par A. $13$\quad B. $18$\quad C. $10$\quad D. $9$
\item Calculate $\frac12+\frac13$.\par A. $\frac{2}{5}$\quad B. $\frac{1}{6}$\quad C. $\frac{5}{3}$\quad D. $\frac{5}{6}$
\item Calculate $\dfrac{2^4}{2^{2}}$.\par A. $2$\quad B. $4$\quad C. $1$\quad D. $9$
\item A value of 50 increases by 20 \%. Final value?\par A. $40$\quad B. $10$\quad C. $60$\quad D. $70$
\item A value of 80 is obtained after a 20 \% decrease. Initial value?\par A. $64$\quad B. $100$\quad C. $96$\quad D. $60$
\item A 25 \% increase then a 20 \% decrease: overall coefficient?\par A. $\frac{9}{10}$\quad B. $\frac{4}{5}$\quad C. $1$\quad D. $\frac{21}{20}$
\item The solution of $2x+4=0$ is\par A. $2$\quad B. $8$\quad C. $-8$\quad D. $-2$
\item The slope between $(0;4)$ and $(2;10)$ is\par A. $-3$\quad B. $3$\quad C. $2$\quad D. $6$
\item The mean of $4, 6, 8$ is\par A. $7$\quad B. $18$\quad C. $6$\quad D. $4$
\item A sector of $90^\circ$ represents what fraction of the total?\par A. $\frac{1}{2}$\quad B. $\frac{1}{3}$\quad C. $\frac{3}{4}$\quad D. $\frac{1}{4}$
\item A probability is $\dfrac{2}{5}$. Its complement is\par A. $\frac{1}{5}$\quad B. $\frac{3}{5}$\quad C. $\frac{2}{5}$\quad D. $\frac{5}{2}$
\item For $f(x)=x^2+4$, $f(-2)$ is\par A. $16$\quad B. $0$\quad C. $6$\quad D. $8$
\end{enumerate}
\mxnewpage\subsection*{Part 2 -- Independent exercises (14 points)}
Every step must be justified. Exact results are expected unless stated otherwise.
\mxneed{12}
\subsection*{Exercise 1 -- An unknown probability (5 points)}
Two events $A$ and $B$ of the same random experiment satisfy $P(A)=\dfrac25$, $P(B)=\dfrac12$ and $P_A(B)=\dfrac34$. We wish to complete the probabilities of a tree whose first level distinguishes $A$ from its complement.
\begin{enumerate}
\item Calculate $P(A\cap B)$. \hfill(1 pt)
\item Determine $P_{\overline A}(B)$ then complete a tree. \hfill(2 pts)
\item Determine $P_B(A)$ and discuss independence. \hfill(2 pts)
\end{enumerate}
\mxneed{12}
\subsection*{Exercise 2 -- Relative position and variation (5 points)}
The functions $f$ and $g$ are defined on $[-1;4]$ by $f(x)=x^3-3x^2+2$ and $g(x)=2$. Their curves are drawn in the same frame. The variation of $f$ must be distinguished from the position of its curve relative to that of $g$.
\begin{enumerate}
\item Calculate $f\prime$ and study the variation. \hfill(1 pt)
\item Compare the curves of f and g on the domain. \hfill(2 pts)
\item Determine the minimum and maximum of f on this domain. \hfill(2 pts)
\end{enumerate}
\mxneed{12}
\subsection*{Exercise 3 -- A perpendicular bisector (4 points)}
In an orthonormal frame, given the points $A(0;2)$, $B(6;4)$ and $P(2;6)$. Let $I$ be the midpoint of $[AB]$. We look for an equation of the perpendicular bisector of this segment, then use this line to compare two distances.
\begin{enumerate}
\item Give the midpoint of AB and a normal to its perpendicular bisector. \hfill(1 pt)
\item Determine the equation of the perpendicular bisector. \hfill(1 pt)
\item Is P equidistant from A and B? Check using two methods. \hfill(2 pts)
\end{enumerate}
